\section{Polynomial Representations}
\label{sec:representations}

This section fixes the notation used throughout the paper and defines the
three univariate representations compared experimentally.

\subsection{Ambient polynomial space}

Let \(\mathbb{F}\) be a field and let
\[
N=2^m
\]
for an integer \(m\ge 1\).  We work in the \(N\)-dimensional
\(\mathbb{F}\)-vector space
\[
\mathbb{F}[X]_{<N}
=
\{P\in\mathbb{F}[X] : \deg P < N\}.
\]
Every representation considered below encodes exactly one polynomial in
this space using \(N\) field elements.

For an integer
\[
e=\sum_{i=0}^{m-1} e_i 2^i,
\qquad e_i\in\{0,1\},
\]
we write
\[
\operatorname{bits}(e)=(e_0,\ldots,e_{m-1}).
\]
For Boolean vectors \(u,v\in\{0,1\}^m\), we write
\[
u\subseteq v
\]
when \(u_i\le v_i\) for every \(i\).  Equivalently, the support of \(u\)
is a subset of the support of \(v\).

\subsection{Monomial representation}

The standard monomial basis of \(\mathbb{F}[X]_{<N}\) is
\[
1,X,X^2,\ldots,X^{N-1}.
\]
A polynomial is represented by coefficients
\[
a_0,\ldots,a_{N-1}\in\mathbb{F}
\]
such that
\[
P(X)=\sum_{e=0}^{N-1}a_eX^e.
\]

For a point \(x\in\mathbb{F}\), Horner evaluation computes
\[
P(x)
=
a_0+x\bigl(a_1+x(a_2+\cdots+x a_{N-1})\bigr),
\]
or equivalently by processing the coefficients in reverse order with the
recurrence
\[
r_{j+1}=xr_j+a_{N-1-j},
\qquad r_0=0.
\]
The serial dependency chain makes ordinary Horner evaluation inherently
sequential, although a long coefficient vector can be partitioned into
blocks and the blocks evaluated in parallel.  Our experimental comparison
therefore includes both serial Horner evaluation and a blocked parallel
monomial baseline.

\subsection{Roots-of-unity Lagrange representation}

Assume that \(\mathbb{F}\) contains a primitive \(N\)-th root of unity
\(\omega\).  Define the multiplicative evaluation domain
\[
H
=
\{1,\omega,\omega^2,\ldots,\omega^{N-1}\}.
\]
A polynomial \(P\in\mathbb{F}[X]_{<N}\) can be represented by its values
\[
v_j=P(\omega^j),
\qquad 0\le j<N.
\]

The associated Lagrange basis consists of polynomials
\[
L_j(X)
=
\prod_{\substack{0\le k<N\\k\ne j}}
\frac{X-\omega^k}{\omega^j-\omega^k},
\]
which satisfy
\[
L_j(\omega^k)=
\begin{cases}
1,&j=k,\\
0,&j\ne k.
\end{cases}
\]
Hence
\[
P(X)=\sum_{j=0}^{N-1}v_jL_j(X).
\]

For a roots-of-unity domain one may also use
\[
L_j(X)
=
\frac{X^N-1}{N(X-\omega^j)}\,\omega^j,
\]
provided \(X\ne\omega^j\).  This identity is useful for arbitrary-point
evaluation because all denominators \(x-\omega^j\) can be inverted with a
single batch inversion.

This representation is especially natural when a proof system already
stores polynomial values on an FFT/NTT domain.  Our benchmark measures
arbitrary-point evaluation from this representation rather than the
separate cost of constructing the representation.

\subsection{Boolean-kernel representation}

For each Boolean vector
\[
y=(y_0,\ldots,y_{m-1})\in\{0,1\}^m,
\]
define
\[
K_y(X)
=
\prod_{i=0}^{m-1}
\left(X^{2^i}+y_i\right).
\]
There are exactly \(2^m=N\) such polynomials.

A Boolean-kernel coefficient vector is a family
\[
(c_y)_{y\in\{0,1\}^m}
\]
representing the polynomial
\[
P(X)
=
\sum_{y\in\{0,1\}^m}
c_yK_y(X).
\]

Although every \(K_y\) has degree at most \(N-1\), and in fact every
\(K_y\) contains the top monomial \(X^{N-1}\), the family is nevertheless
a basis.  Its coefficient matrix is a Boolean zeta matrix after a reversal
of the monomial index.

\begin{lemma}[Monomial coefficients of \(K_y\)]
\label{lem:kernel-monomial-coeff}
Let
\[
e=\sum_{i=0}^{m-1}e_i2^i
\]
with \(e_i\in\{0,1\}\), and define the Boolean complement vector
\[
\bar e=(1-e_0,\ldots,1-e_{m-1}).
\]
Then
\[
[X^e]K_y(X)
=
\begin{cases}
1,&\bar e\subseteq y,\\
0,&\text{otherwise}.
\end{cases}
\]
\end{lemma}

\begin{proof}
Expanding
\[
K_y(X)=\prod_{i=0}^{m-1}(X^{2^i}+y_i),
\]
the monomial \(X^e\) can arise in exactly one way because the powers
\(2^0,\ldots,2^{m-1}\) give the unique binary expansion of \(e\).

For each index \(i\) with \(e_i=1\), the term \(X^{2^i}\) must be selected.
For each index \(i\) with \(e_i=0\), the constant term \(y_i\) must be
selected.  Therefore
\[
[X^e]K_y(X)
=
\prod_{i:e_i=0}y_i.
\]
Since every \(y_i\) is either \(0\) or \(1\), this product equals \(1\)
exactly when \(y_i=1\) for every \(i\) such that \(e_i=0\).  This condition
is precisely
\[
\bar e\subseteq y.
\]
\end{proof}

\begin{theorem}[Boolean-kernel basis]
\label{thm:kernel-basis}
The family
\[
\mathcal{K}_m
=
\{K_y(X):y\in\{0,1\}^m\}
\]
is a basis of \(\mathbb{F}[X]_{<N}\).
\end{theorem}

\begin{proof}
There are exactly \(N\) polynomials in \(\mathcal{K}_m\), so it is enough
to prove linear independence.

Index the monomial rows not directly by \(e\), but by the Boolean vector
\[
z=\bar e.
\]
By \cref{lem:kernel-monomial-coeff}, the coefficient matrix \(A\) of the
family \(\mathcal{K}_m\) in the monomial basis is
\[
A_{z,y}
=
\mathbf{1}[z\subseteq y].
\]
This is the zeta matrix of the Boolean subset lattice.

Order the Boolean vectors by nondecreasing Hamming weight, breaking ties
arbitrarily.  Whenever \(z\subseteq y\), either \(z=y\) or
\[
|z|<|y|.
\]
Consequently, under this ordering, \(A\) is upper triangular with diagonal
entries
\[
A_{y,y}=1.
\]
Hence
\[
\det A=1\ne0
\]
in every field, so \(A\) is invertible.  Therefore the polynomials
\(K_y\) are linearly independent and form a basis of
\(\mathbb{F}[X]_{<N}\).
\end{proof}

\subsection{Relation to the Boolean zeta transform}

The proof above also gives an explicit relationship between monomial and
Boolean-kernel coefficients.

Let
\[
P(X)=\sum_{e=0}^{N-1}a_eX^e
=
\sum_{y\in\{0,1\}^m}c_yK_y(X).
\]
For the exponent whose Boolean complement is \(z\), that is
\[
e=(N-1)\mathbin{\mathrm{xor}} z,
\]
\cref{lem:kernel-monomial-coeff} gives
\[
a_e
=
\sum_{y:\,z\subseteq y}c_y.
\]
Thus, after complementing the monomial index, the map from kernel
coefficients to monomial coefficients is the Boolean superset-zeta
transform.

The inverse map is Boolean Möbius inversion:
\[
c_y
=
\sum_{z:\,y\subseteq z}
(-1)^{|z|-|y|}
a_{(N-1)\mathbin{\mathrm{xor}} z}.
\]
Algorithmically, both transforms can be implemented in
\[
O(N\log N)
\]
field operations by the usual butterfly-style zeta/Möbius transform.

This conversion cost is deliberately separated from the single-point
evaluation benchmarks.  In each benchmark the polynomial is assumed to
already be available in the representation being measured.  Conversion
costs become important when considering an end-to-end proof-system
pipeline and are discussed separately later in the paper.

\subsection{A small example}

Take \(m=2\), so \(N=4\).  The four kernel polynomials are
\[
\begin{aligned}
K_{00}(X)&=X(X^2)=X^3,\\
K_{01}(X)&=(X+1)X^2=X^3+X^2,\\
K_{10}(X)&=X(X^2+1)=X^3+X,\\
K_{11}(X)&=(X+1)(X^2+1)=X^3+X^2+X+1.
\end{aligned}
\]
Here the first Boolean coordinate corresponds to the factor
\(X^{2^0}+y_0\), and the second to \(X^{2^1}+y_1\).

The coefficient pattern is triangular after complementing the monomial
index.  In particular, the constant monomial occurs only in \(K_{11}\),
the monomial \(X\) occurs in \(K_{01}\) and \(K_{11}\), the monomial
\(X^2\) occurs in \(K_{10}\) and \(K_{11}\), and \(X^3\) occurs in all
four kernel polynomials.

This subset structure is the algebraic reason that the representation is
invertible.  In the next section we show that the same product structure
also produces a particularly simple binary evaluation algorithm.
